
\documentclass[11pt]{article}

% ------------------------------------------------------------
% Packages
% ------------------------------------------------------------
\usepackage[margin=1in]{geometry}
\usepackage[T1]{fontenc}
\usepackage{lmodern}
\usepackage{pdflscape}
\usepackage{longtable}
\usepackage{microtype}
\usepackage{amsmath,amssymb,mathtools}
\usepackage{booktabs}
\usepackage{array}
\usepackage{graphicx}
\usepackage{caption}
\usepackage{subcaption}
\usepackage{enumitem}
\usepackage{xcolor}
\usepackage{hyperref}
\usepackage{setspace}
\usepackage{natbib}
\usepackage{threeparttable}

\hypersetup{
    colorlinks=true,
    linkcolor=blue,
    citecolor=blue,
    urlcolor=blue
}

\setlength{\parindent}{0pt}
\setlength{\parskip}{0.72em}
\onehalfspacing


% ------------------------------------------------------------
% Title information
% ------------------------------------------------------------
\title{\textbf{Sex Differences in the Association Between Insulin Resistance and Systemic Inflammation}\\
\textbf{Below the HbA1c Prediabetes Threshold: NHANES 2015--March 2020}}

% ------------------------------------------------------------
% Document
% ------------------------------------------------------------
\begin{document}

\author{
Scott Moerschbacher$^{1}$,
Janice Baker$^{1,2}$
}

\date{}
\maketitle

\begin{center}
$^{1}$Open Inquiry Research Group, USA\\
$^{2}$Department of Sociology, University of Nevada, Reno,
Reno, Nevada, USA
\end{center}

\vspace{0.75em}

\begin{center}
\textbf{Corresponding author:}\\
Scott Moerschbacher\\
Open Inquiry Research Group\\
\href{mailto:scott@openinquiryresearch.org}
{scott@openinquiryresearch.org}
\end{center}

\vspace{0.5em}

\noindent\textbf{Keywords:}
insulin resistance; HOMA-IR; systemic inflammation; hsCRP;
sex differences; NHANES; glycemic threshold

\vspace{1em}

\begin{abstract}

\textbf{Background:}
Insulin resistance and systemic inflammation are closely related features of
metabolic dysfunction, but whether their association differs by sex before
conventional glycemic thresholds for prediabetes are reached is not well
established. We examined whether sex modifies the association between
homeostatic model assessment of insulin resistance (HOMA-IR) and
high-sensitivity C-reactive protein (hsCRP) among U.S. adults without
diagnosed diabetes and with HbA1c $<5.7\%$.

\textbf{Methods:}
We conducted a cross-sectional analysis of the 2015--March 2020
pre-pandemic National Health and Nutrition Examination Survey (NHANES).
The primary analytic sample included 3,079 adults aged $\geq20$ years
meeting the glycemic and fasting-subsample eligibility criteria.
Survey-weighted linear regression modeled log-transformed hsCRP as a
function of HOMA-IR, sex, their interaction, and period-specific adjustment
for HbA1c, waist circumference, age, and non-HDL cholesterol. The primary
inferential test was the survey-adjusted HOMA-IR-by-sex interaction.
Sensitivity analyses evaluated exposure scale and upper-tail influence,
alternative adiposity specifications, selected clinical and behavioral
factors, pregnancy, and restriction to hsCRP $\leq10$ mg/L.

\textbf{Results:}
The adjusted HOMA-IR--hsCRP association differed by sex in the primary
analysis ($F_{1,40}=9.05$, $p=0.0045$). The estimated HOMA-IR slope was
0.0629 (95\% CI, 0.0265 to 0.0992) among women and 0.0082
(95\% CI, $-0.0156$ to 0.0319) among men, with a male-minus-female slope
difference of $-0.0547$ (95\% CI, $-0.0914$ to $-0.0180$).
The direction of the sex contrast was maintained across alternative HOMA-IR
scales and upper-tail restrictions, alternative adiposity specifications,
and targeted smoking, statin, renal, hepatic, alcohol, and pregnancy
sensitivities. However, restriction to hsCRP $\leq10$ mg/L attenuated the
male-minus-female slope difference by approximately 51\%, to $-0.0269$
(95\% CI, $-0.0572$ to 0.0033; $F_{1,40}=3.24$, $p=0.080$).

\textbf{Conclusions:}
Among U.S. adults without diagnosed diabetes and with HbA1c below the
conventional threshold for prediabetes, the adjusted association between
HOMA-IR and hsCRP was more positive among females than among males. This sex contrast
was stable across multiple metabolic and clinical sensitivity analyses but
was substantially attenuated after restriction of the upper hsCRP range.
These findings suggest that sex-related heterogeneity in
metabolic--inflammatory coupling may be detectable below conventional
glycemic thresholds, while also indicating that the magnitude of this
heterogeneity depends importantly on the inflammatory range represented in
the analytic population.

\end{abstract}

% ============================================================

\section{Introduction}

Insulin resistance and systemic inflammation are closely related features
of metabolic dysfunction. Higher concentrations of C-reactive protein (CRP)
have been associated with greater insulin resistance in population-based
studies of adults without diabetes, including nationally representative
U.S. data \citep{Meng2007CRPInsulinResistance}. Longitudinal evidence from
the Coronary Artery Risk Development in Young Adults (CARDIA) study has
further demonstrated associations between CRP and both concurrent and
subsequent HOMA-IR, although much of the association was attenuated after
accounting for adiposity \citep{Park2009InflammationHOMA}. Similar
associations between CRP and HOMA-IR have been reported in other
nondiabetic populations and have persisted after adjustment for abdominal
adiposity \citep{Chou2010CRPInsulinResistance}. These observations suggest
that inflammatory and insulin-resistance phenotypes may be coupled before
the development of overt diabetes.

Sex may be an important source of heterogeneity in this relationship. Women
and men differ in insulin sensitivity and in the distribution and metabolic
function of adipose tissue, with sex hormones and inflammatory pathways
among the mechanisms proposed to contribute to these differences
\citep{Gado2024SexIR}. Immune and inflammatory responses also differ by sex
across both innate and adaptive pathways
\citep{KleinFlanagan2016SexImmune}. Consistent with these physiological
differences, women have generally been reported to have higher CRP
concentrations than men in population-based studies
\citep{Lakoski2006GenderCRP}, and the association between total fat mass and
CRP may itself differ by sex \citep{Khera2009SexCRPBodyFat}. Thus, sex
differences in average concentrations of metabolic or inflammatory markers
may not fully characterize sex-related heterogeneity in the relationship
between these systems.

Previous studies have examined the relationship between HOMA-IR and CRP in
women and men, but evidence regarding sex-specific coupling is limited and
not uniform. Sex-stratified analyses in nondiabetic Taiwanese adults found
positive adjusted HOMA-IR--CRP associations in both women and men
\citep{Chou2010CRPInsulinResistance}, whereas CARDIA reported no significant
interaction between CRP and sex in the prediction of HOMA-IR
\citep{Park2009InflammationHOMA}. More recently, analysis of a representative
Canadian adult population explicitly examined potential differences by sex
and reported differing patterns of association between CRP and HOMA-IR in
women and men \citep{Shahid2023CRPHOMA}. Less is known about whether the
adjusted association between insulin resistance and systemic inflammation
differs by sex among adults who remain below conventional glycemic thresholds
for prediabetes. This distinction may be informative because clinical
glycemic categories do not necessarily imply physiological homogeneity in
insulin resistance or inflammatory state.

We therefore examined whether sex modifies the association between HOMA-IR
and high-sensitivity C-reactive protein (hsCRP) among U.S. adults aged 20
years or older without diagnosed diabetes and with HbA1c $<5.7\%$, using
data from the 2015--March 2020 pre-pandemic National Health and Nutrition
Examination Survey (NHANES). The primary analysis formally tested the
HOMA-IR-by-sex interaction while accounting for the complex survey design
and relevant metabolic covariates. We further evaluated the stability of
the estimated sex difference across alternative HOMA-IR scales and exposure
distributions, multiple representations of adiposity, selected clinical and
behavioral factors, pregnancy, and restriction of the upper hsCRP range.

% ============================================================

\section{Methods}

\subsection{Study Design and Data Source}

We conducted a cross-sectional analysis of publicly available data from the
National Health and Nutrition Examination Survey (NHANES), using the
2015--2016 cycle and the 2017--March 2020 pre-pandemic data release.
NHANES is a nationally representative survey of the noninstitutionalized
U.S. population that combines interviews, physical examinations, and
laboratory measurements using a complex, multistage probability sampling
design \citep{NHANES2017_2020Guidelines}.

The 2015--2016 and 2017--March 2020 datasets were harmonized at the
participant level using the NHANES participant identifier (\texttt{SEQN}).
Variables required for the primary analysis were obtained from the
demographic, examination, diabetes questionnaire, glycohemoglobin,
fasting glucose, fasting insulin, high-sensitivity C-reactive protein,
and lipid files. Additional questionnaire and laboratory files were used for sensitivity analyses. 
Because HOMA-IR requires both fasting glucose and fasting insulin, analyses were conducted within the
NHANES fasting laboratory subsample and used the corresponding
fasting-subsample survey weights.

\subsection{Analytic Population}

The analytic population was restricted to adults aged 20 years or older
who had no self-reported diagnosis of diabetes and had HbA1c $<5.7\%$,
the lower boundary of the conventional HbA1c-defined prediabetes range
\citep{ADA2026Diagnosis}. Participants were additionally required to belong
to the qualifying NHANES fasting laboratory subsample, as indicated by a
nonzero fasting-subsample survey weight, and to have valid fasting glucose
and fasting insulin measurements, positive HOMA-IR and hsCRP values, and
complete data for variables included in the primary model. No imputation
of missing data was performed.

Fasting-subsample eligibility was independently verified using
participant-reported fasting duration from the NHANES fasting questionnaire,
calculated as \texttt{PHAFSTHR} + \texttt{PHAFSTMN}/60
\citep{NHANESFastingQuestionnaire2015,NHANESFastingQuestionnaire2017_2020}. All participants with a
qualifying fasting-subsample weight in the final analytic population reported
fasting durations within the NHANES qualifying range of 8 to $<24$ hours.

Of 3,182 participants meeting the remaining eligibility and measurement
requirements, 103 (3.2\%) were excluded because of incomplete primary-model
covariate data, yielding a final primary analytic sample of 3,079 participants.
The final sample represented approximately 144.3 million U.S. adults after
application of the combined survey weights.

No a priori sample-size calculation was performed; the analytic sample
comprised all participants from the specified NHANES survey periods who met
the study eligibility criteria and had the data required for the primary
analysis.

The HbA1c restriction was used to examine metabolic heterogeneity specifically
among individuals remaining below the conventional laboratory threshold for
prediabetes. This restriction defines the population under study and should
not be interpreted as establishing that all included participants were
otherwise metabolically healthy.

\subsection{Exposure: HOMA-IR}

The primary exposure was the Homeostatic Model Assessment of Insulin
Resistance (HOMA-IR), a surrogate measure of insulin resistance derived from
fasting insulin and fasting glucose concentrations
\citep{Matthews1985HOMA,Wallace2004HOMA}. HOMA-IR was calculated as

\[
\mathrm{HOMA\!-\!IR}
=
\frac{
\mathrm{fasting\ insulin}\;(\mu\mathrm{U/mL})
\times
\mathrm{fasting\ glucose}\;(\mathrm{mg/dL})
}{405}.
\]

HOMA-IR was modeled continuously and retained on its original scale in the
primary analysis. Raw HOMA-IR was retained to preserve direct interpretation
per one-unit difference and continuity with the preceding analysis of this
NHANES cohort. Because HOMA-IR was right-skewed, natural-log-transformed
HOMA-IR and progressively trimmed upper-tail samples were evaluated in
sensitivity analyses.

\subsection{Outcome: hsCRP}

The primary outcome was high-sensitivity C-reactive protein (hsCRP), modeled
as its natural logarithm,

\[
Y=\ln(\mathrm{hsCRP}).
\]

The log transformation was retained from the preceding analysis of this
cohort because raw hsCRP was strongly right-skewed and log transformation
substantially improved regression residual behavior. For models using
$\ln(\mathrm{hsCRP})$, coefficients may be expressed as proportional
differences on the original hsCRP scale according to

\[
100\left(e^\beta-1\right)\%.
\]

A sensitivity analysis specified before examination of the corresponding
Study \#2 result restricted the analytic population to participants with
hsCRP $\leq10$ mg/L to evaluate dependence of the estimated
sex interaction on the upper range of the hsCRP distribution.

\subsection{Sex and Primary Covariates}

Sex was the focal effect modifier and was coded from the NHANES sex variable
as female (reference category) or male. The primary analysis therefore tested
whether the adjusted HOMA-IR--hsCRP slope differed between females and males;
sex was not included as a separate ordinary adjustment covariate outside the
interaction term.

The primary harmonized adjustment set consisted of HbA1c, waist circumference,
continuous age, and non-HDL cholesterol. Non-HDL cholesterol was calculated as

\[
\mathrm{non\!-\!HDL}
=
\mathrm{total\ cholesterol}
-
\mathrm{HDL\ cholesterol}.
\]

Waist circumference was retained as the primary adiposity measure. Alternative
adiposity specifications using body mass index (BMI) and waist-to-height ratio,
as well as a model without adiposity adjustment, were evaluated as sensitivity
analyses.

\subsection{NHANES Survey Design and Weighting}

All population estimates and inferential analyses incorporated the NHANES
complex sampling design \citep{NHANES2017_2020Guidelines}. Because calculation
of HOMA-IR required both fasting glucose and fasting insulin, analyses used the
corresponding fasting-subsample survey weights rather than general examination
weights.

Survey variance estimation incorporated the NHANES masked variance stratum
(\texttt{SDMVSTRA}) and masked variance unit (\texttt{SDMVPSU}) variables.
For the combined 2015--March 2020 analysis, fasting weights were adjusted for
the unequal durations represented by the component datasets according to
NHANES guidance \citep{NHANES2017_2020Guidelines}. The 2015--2016 fasting
weight was rescaled as

\[
W_{15/16}^{*}
=
\mathrm{WTSAF2YR}
\left(\frac{2}{5.2}\right),
\]

and the 2017--March 2020 pre-pandemic fasting weight as

\[
W_{17/20}^{*}
=
\mathrm{WTSAFPRP}
\left(\frac{3.2}{5.2}\right).
\]

Masked variance strata and variance units were uniquely identified across
periods before the datasets were combined. The resulting combined design
provided 40 survey design degrees of freedom. Survey-design-based variance
estimates were used for standard errors, confidence intervals, $t$-tests,
and $F$-tests; participants were not treated as independent
simple-random-sample observations.

Analyses were performed on the full survey design while defining the eligible
analytic population as a survey domain. Contributions to estimating equations
from observations outside the analytic domain were set to zero, preserving the
design structure for domain variance estimation.

\subsection{Primary Statistical Analysis}

The primary analysis estimated the adjusted association between HOMA-IR and
log-transformed hsCRP and formally tested whether that association differed
by sex. The Study \#2 primary interaction model was

\[
\ln(\mathrm{hsCRP})
\sim
\mathrm{HOMA\!-\!IR}\times\mathrm{sex}
+
\mathrm{period}\times
(\mathrm{HbA1c}+\mathrm{waist}+\mathrm{age}+\mathrm{non\!-\!HDL}).
\]

Equivalently, the focal portion of the model may be represented as

\[
\ln(\mathrm{hsCRP}_i)
=
\beta_0
+\beta_1 H_i
+\beta_2 M_i
+\beta_3(H_iM_i)
+\text{period-specific adjustment terms}
+\epsilon_i,
\]

where $H_i$ denotes HOMA-IR and $M_i$ is an indicator for male sex. With
female as the reference category, $\beta_1$ represents the adjusted female
HOMA-IR slope and $\beta_3$ represents the male-minus-female difference in
HOMA-IR slopes. The male slope is therefore $\beta_1+\beta_3$.

The primary inferential test was the survey-adjusted Wald test of

\[
H_0:\beta_3=0.
\]

Evidence for effect modification was therefore determined from the formal
HOMA-IR-by-sex interaction rather than by comparing whether the female and
male sex-specific coefficients separately crossed a conventional statistical
significance threshold. Sex-specific slopes, the male-minus-female slope
contrast, and corresponding 95\% confidence intervals were estimated from the
same fitted interaction model.

A survey-period indicator was included, and HbA1c, waist circumference,
continuous age, and non-HDL cholesterol were permitted to have period-specific
coefficients through interactions with survey period. The HOMA-IR-by-sex
association remained the common effect-modification structure of primary
interest. This specification avoided assuming that adjustment-variable
relationships with hsCRP were identical in the 2015--2016 and
2017--March 2020 components.

\subsection{Sensitivity Analyses}

Sensitivity analyses were organized according to specific potential sources of
inferential vulnerability and were treated as robustness analyses rather than
alternative primary models.

\paragraph{Exposure scale and upper-tail influence.}
To assess whether the sex interaction depended on the right-skewed raw
HOMA-IR scale or a small number of extreme exposure values, the primary model
was repeated using natural-log-transformed HOMA-IR and after excluding
observations above the 99th, 97.5th, and 95th percentiles of the HOMA-IR
distribution. Residual diagnostics for the raw- and log-HOMA-IR
specifications were also examined.

\paragraph{Adiposity specification.}
The primary waist-circumference model was compared with specifications using
BMI, waist-to-height ratio, and no adiposity adjustment. Both available-case
and common-complete-case analyses were used to distinguish changes associated
with adiposity specification from those caused by differential missingness.

\paragraph{Additional adjustment sensitivities.}
Targeted sensitivity analyses evaluated additional adjustment for smoking
status, statin use, renal function, hepatic/metabolic
biochemistry, and alcohol-use status. Smoking was categorized as never,
former, or current. Statin use was identified from reported prescription
medications during the preceding 30 days using NHANES drug classifications for
HMG-CoA reductase inhibitors. Renal sensitivity analyses used continuous
2021 CKD-EPI creatinine-based estimated glomerular filtration rate (eGFR)
\citep{Inker2021CKDEPI}, with log-transformed serum creatinine as a
complementary specification. Log-transformed alanine aminotransferase (ALT) was used for the
hepatic/metabolic sensitivity. Alcohol use was reconstructed across the two
questionnaire structures as zero past-year exposure, very-low lifetime
exposure, or positive past-year drinking frequency. For each sensitivity with
missing exposure information, the reference and expanded models were compared
within the same common sample.

\paragraph{Current pregnancy.}
Because current pregnancy represents a physiologically distinct metabolic and
inflammatory state, the primary model was repeated after excluding participants
with confirmed current pregnancy. A complementary analysis additionally
excluded participants whose pregnancy status could not be determined. These
analyses addressed current pregnancy only and were not interpreted as
accounting for menopausal or other reproductive states.

\paragraph{Upper hsCRP restriction.}
The primary model was repeated after restriction to hsCRP $\leq10$ mg/L.
Elevations above 10 mg/L have historically received particular attention in
the interpretation of CRP, although such elevations need not represent acute
inflammation and may be repeatedly observed in some individuals
\citep{Ishii2012CRP}. This threshold had been used in the preceding analysis
of the same cohort and was not selected from the Study \#2 interaction results.
No alternative hsCRP thresholds were explored after observing this sensitivity
analysis.

\subsection{Descriptive Follow-up of the Upper hsCRP Range}

Because restriction to hsCRP $\leq10$ mg/L materially attenuated the
estimated HOMA-IR-by-sex interaction, a descriptive follow-up was conducted to
characterize participants with hsCRP $>10$ mg/L. This analysis was undertaken
to understand the sensitivity result and was not used to define new exclusion
criteria or alternative primary models.

A limited complete-blood-count panel was specified before examining its values:
total white blood cell count, absolute neutrophil count, neutrophil percentage,
absolute lymphocyte count, absolute monocyte count, and a derived
neutrophil-to-lymphocyte ratio. These measures were summarized descriptively by
sex and hsCRP group. The HOMA-IR distribution was likewise summarized within
the four sex-by-hsCRP groups. Scatterplots of raw HOMA-IR against log-hsCRP
were examined separately by sex using both the full HOMA-IR range and a
HOMA-IR 0--12 display window; the latter was a visualization restriction only
and did not exclude observations from any analysis.

The descriptive CBC analysis was not used to classify participants as having
acute infection. Similarly, the joint HOMA-IR/hsCRP patterns were interpreted
as descriptive features of the observed distribution rather than as evidence
of a causal biological mechanism.

\subsection{Deferred Reproductive-State Analysis}

Reproductive-health questionnaire variables describing recent menses,
hysterectomy, and bilateral oophorectomy were examined as predictor-only
quality-control work after the sex interaction had been identified. The
available cross-sectional information did not provide a uniquely determined
menopausal classification for all women, and no HOMA-IR--hsCRP outcome model
was fitted using post hoc reproductive-state categories. Reproductive and
menopausal physiology was therefore deferred as a separate future research
question rather than added as another Study \#2 sensitivity analysis.

\subsection{Software and Reproducibility}

Analyses were performed in Python using NumPy, pandas, Patsy, SciPy, and
statsmodels. Survey-weighted linear regression used Taylor-linearized
design-based covariance estimation with NHANES strata, primary sampling units,
fasting-subsample weights, and domain score-zeroing. The validated survey
analysis implementation and cohort-construction framework from the preceding
NHANES study were retained without modification for the present analysis.

All sensitivity analyses were documented sequentially in an analysis log,
including the specification of exposure definitions and analytic decisions
before outcome-model fitting where applicable. Analyses were frozen after
completion; additional thresholds or model variants were not selected to
recover statistical significance.

\subsection{Use of Generative Artificial Intelligence}

Generative artificial intelligence (ChatGPT, OpenAI) was used as a
research-support tool during this project, including assistance with drafting
and revising portions of the manuscript and with writing, optimizing, and
debugging portions of the reproducible Python analysis pipeline. The authors
determined the research questions, analytic strategy, statistical methods,
interpretation of results, and conclusions. AI-assisted analytical code and
outputs were subject to human review and validation against the specified
analytic framework, mathematical expectations, and NHANES survey-design
constraints. All manuscript content was critically reviewed and approved by
the authors, who take full responsibility for the accuracy, integrity,
interpretation, and conclusions of the work.

% ============================================================

\section{Results}

\subsection{Analytic Population}

The final primary analytic sample included 3,079 adults aged 20 years or older
without diagnosed diabetes and with HbA1c $<5.7\%$. After application of the
combined NHANES fasting-subsample weights, this sample represented
approximately 144.3 million U.S. adults. The complex survey design provided
40 denominator degrees of freedom for design-based inference. Sex-stratified descriptive characteristics are shown in Table~\ref{tab:table1}.


\begin{table}[htbp]
\centering
\caption{Characteristics of the primary analytic population by sex}
\label{tab:table1}
\begin{threeparttable}
\begin{tabular}{lcc}
\toprule
Characteristic & Female ($n=1,638$) & Male ($n=1,441$) \\
\midrule
Weighted population, millions & 74.2 & 70.2 \\
Age, years & 43.6 (0.6) & 42.0 (0.6) \\
HbA1c, \% & 5.26 (0.01) & 5.25 (0.01) \\
Waist circumference, cm & 94.4 (0.6) & 99.2 (0.6) \\
Non-HDL cholesterol, mg/dL & 126.6 (1.4) & 135.1 (1.6) \\
Fasting glucose, mg/dL & 97.3 (0.3) & 101.5 (0.4) \\
Fasting insulin, $\mu$U/mL & 7.76 (5.25--12.01) & 7.86 (5.01--13.74) \\
HOMA-IR & 1.87 (1.21--2.88) & 2.00 (1.25--3.43) \\
hsCRP, mg/L & 1.70 (0.67--3.93) & 1.30 (0.61--2.67) \\
\bottomrule
\end{tabular}
\begin{tablenotes}[flushleft]
\footnotesize
\item Values are survey-weighted mean (SE) for age, HbA1c, waist circumference,
non-HDL cholesterol, and fasting glucose, and survey-weighted median (IQR) for
fasting insulin, HOMA-IR, and hsCRP.
\item The analytic population included adults aged $\geq20$ years without
diagnosed diabetes and with HbA1c $<5.7\%$, meeting the fasting-subsample and
complete-case criteria for the primary model.
\item Weighted population estimates use the combined NHANES 2015--2016 and
2017--March 2020 pre-pandemic fasting-subsample weights.
\item No hypothesis tests comparing baseline characteristics by sex were performed.
\end{tablenotes}
\end{threeparttable}
\end{table}


\subsection{Primary HOMA-IR-by-Sex Interaction}

In the locked primary model, the adjusted association between HOMA-IR and
log-transformed hsCRP differed by sex
($F(1,40)=9.05$, $p=0.0045$; Table~\ref{tab:primary_sex_interaction}). The estimated HOMA-IR slope among
females was 0.0629 (SE 0.0180; 95\% CI 0.0265 to 0.0992), whereas the
estimated slope among males was 0.0082 (SE 0.0117; 95\% CI $-0.0156$ to
0.0319). The male-minus-female difference in slopes was $-0.0547$
(95\% CI $-0.0914$ to $-0.0180$).

On the original hsCRP scale, these coefficients correspond descriptively to
approximately a 6.5\% increase in hsCRP per one-unit higher HOMA-IR among
females and a 0.8\% increase among males, conditional on the fitted model.
Inference, however, was based on the log-hsCRP scale and on the formal
HOMA-IR-by-sex interaction rather than on separate tests of the two
sex-specific slopes.

\begin{table}[htbp]
\centering
\caption{Primary survey-weighted HOMA-IR-by-sex interaction model}
\label{tab:primary_sex_interaction}
\begin{threeparttable}
\begin{tabular}{lrrrr}
\toprule
Estimate & $\beta$ & SE & 95\% CI & $p$ \\
\midrule
Female HOMA-IR slope & 0.0629 & 0.0180 & 0.0265 to 0.0992 & 0.0012 \\
Male HOMA-IR slope & 0.0082 & 0.0117 & $-0.0156$ to 0.0319 & 0.4905 \\
Male-minus-female slope contrast & $-0.0547$ & --- & $-0.0914$ to $-0.0180$ & 0.0045 \\
\bottomrule
\end{tabular}
\begin{tablenotes}[flushleft]
\footnotesize
\item Outcome: $\ln(\mathrm{hsCRP})$; exposure: raw continuous HOMA-IR.
The model included HOMA-IR$\times$sex and period-specific adjustment for
continuous HbA1c, waist circumference, age, and non-HDL cholesterol.
Female was the reference sex. The formal interaction test for equality of the
sex-specific HOMA-IR slopes was $F(1,40)=9.053$, $p=0.00452$.
Inference incorporated the NHANES fasting-subsample weights, masked strata,
and masked PSUs. Analytic $n=3,079$.
\end{tablenotes}
\end{threeparttable}
\end{table}


\subsection{Exposure Scale and Upper-Tail Influence}

The HOMA-IR distribution was right-skewed (median 1.99; 95th percentile 6.66;
97.5th percentile 8.44; 99th percentile 11.31; maximum 38.97). The estimated
sex contrast persisted in the same direction after natural-log transformation
of HOMA-IR and after progressive restriction of the upper HOMA-IR tail
(Table~S1). In the full sample, the log-HOMA-IR interaction was
$F(1,40)=10.88$ ($p=0.0021$). Under raw-HOMA-IR restrictions at the 99th,
97.5th, and 95th percentiles, male-minus-female slope contrasts were
$-0.0735$, $-0.1036$, and $-0.1001$, respectively, compared with $-0.0547$
in the full primary sample. Residual diagnostics did not identify a
qualitative deficiency requiring replacement of the primary raw-HOMA-IR
specification.

\subsection{Adiposity Sensitivity}

Adiposity adjustment substantially attenuated the HOMA-IR--hsCRP slopes
relative to a model without adiposity adjustment, but the estimated sex
contrast remained in the same direction across BMI, waist circumference, and
waist-to-height ratio specifications (Table~S2). The male-minus-female
contrast was $-0.0955$ (95\% CI $-0.1616$ to $-0.0293$) without adiposity
adjustment, $-0.0423$ (95\% CI $-0.0811$ to $-0.0035$) with BMI,
$-0.0547$ (95\% CI $-0.0914$ to $-0.0180$) with the primary waist
circumference specification, and $-0.0416$ (95\% CI $-0.0791$ to
$-0.0041$) with waist-to-height ratio. Common-complete-case analyses
($n=3,072$) produced nearly identical estimates, indicating that differences
among these specifications were not attributable to differential missingness.

\subsection{Additional Adjustment and Exclusion Sensitivities}

Targeted adjustment for smoking, current prescription statin use, renal
function, log-transformed ALT, and alcohol-use status produced little change
in the estimated HOMA-IR-by-sex contrast within their respective common
samples (Table~S3). The alcohol analysis was particularly informative about
sample restriction: among the 2,919 participants with classifiable alcohol
status, the reference contrast was $-0.0673$ (95\% CI $-0.1060$ to
$-0.0287$), compared with $-0.0665$ (95\% CI $-0.1047$ to $-0.0284$)
after alcohol adjustment. Thus, the difference between the full-cohort
primary estimate and the alcohol-sample estimate reflected sample restriction
rather than alcohol adjustment.

Exclusion of 48 participants with confirmed current pregnancy likewise
produced little change ($n=3,031$; male-minus-female contrast $-0.0567$;
$F(1,40)=9.04$, $p=0.0045$). Excluding both confirmed pregnancy and the
15 participants whose pregnancy status could not be determined produced a
nearly identical contrast ($-0.0567$; $n=3,016$).

\subsection{Sensitivity to the Upper hsCRP Range}

Restriction to participants with hsCRP $\leq10$ mg/L materially attenuated
the estimated sex difference. In the full primary sample, the
male-minus-female HOMA-IR slope contrast was $-0.0547$ (95\% CI $-0.0914$
to $-0.0180$). Among the 2,895 participants with hsCRP $\leq10$ mg/L, the
female slope decreased from 0.0629 to 0.0438, whereas the male slope increased
from 0.0082 to 0.0168. The resulting male-minus-female contrast was
$-0.0269$ (95\% CI $-0.0572$ to 0.0033), with
$F(1,40)=3.24$ ($p=0.0796$).

Thus, restriction to hsCRP $\leq10$ mg/L reduced the absolute magnitude of
the interaction contrast by approximately 50.8\%. The attenuation therefore
reflected convergence of both sex-specific slopes rather than only a change
in whether the interaction crossed a conventional statistical-significance
threshold.

\subsection{Descriptive Characterization of hsCRP $>10$ mg/L}

Of the 184 participants with hsCRP $>10$ mg/L, 128 were female and 56 were
male, corresponding to unweighted proportions of 7.81\% and 3.89\% within
the female and male analytic samples, respectively. Complete-blood-count
measures were available for all 184 participants.

In both sexes, the hsCRP $>10$ mg/L group showed descriptively higher WBC,
absolute neutrophil count, neutrophil percentage, and
neutrophil-to-lymphocyte ratio than the hsCRP $\leq10$ mg/L group
(Table~S4). The HOMA-IR distribution showed a more pronounced shift among
females: median HOMA-IR was 3.01 versus 1.91 for hsCRP $>10$ versus
$\leq10$ mg/L, with corresponding 90th percentiles of 7.40 and 4.51.
Among males, median HOMA-IR was similar between hsCRP groups (2.04 versus
2.05), although the upper distribution extended farther among those with
hsCRP $>10$ mg/L (90th percentile 6.77 versus 5.13).

Sex-stratified scatterplots were consistent with these descriptive summaries:
female observations with hsCRP $>10$ mg/L extended farther into the
higher-HOMA-IR portion of the joint distribution, whereas numerous male
observations with higher HOMA-IR remained below the 10 mg/L hsCRP threshold.
These descriptive analyses were not used to classify acute infection or to
define additional exclusions.


% ============================================================

\section{Discussion}

In this nationally representative sample of U.S. adults without diagnosed
diabetes and with HbA1c below the conventional threshold for prediabetes,
the adjusted association between HOMA-IR and hsCRP differed by sex. The
association was more positive among women than men, with a male-minus-female
difference in HOMA-IR slope of $-0.0547$ (95\% CI, $-0.0914$ to $-0.0180$).
This contrast persisted across alternative HOMA-IR scales, progressively
trimmed exposure distributions, alternative representations of adiposity,
and multiple targeted clinical and behavioral sensitivity analyses. The
principal qualification was the restriction to hsCRP $\leq10$ mg/L, which
reduced the estimated sex contrast by approximately 51\%. Thus, the findings
support sex-related heterogeneity in the relationship between fasting
insulin resistance and systemic inflammation in this metabolically selected
population, while indicating that the magnitude of this heterogeneity is
substantially influenced by the upper range of the hsCRP distribution.

The observed sex difference is biologically plausible in the context of
established sexual dimorphism in metabolic and immune physiology. Sex
differences have been described in insulin sensitivity and resistance, with
potential contributions from body-fat distribution and function,
inflammatory biology, and sex hormones \citep{Gado2024SexIR}. Immune
responses likewise differ by sex across both innate and adaptive pathways
and vary across the life course and reproductive state
\citep{KleinFlanagan2016SexImmune}. At the population level, women have been
reported to have higher CRP concentrations than men even after accounting
for several relevant covariates \citep{Lakoski2006GenderCRP}, and the
association between total fat mass and log-transformed CRP has been reported
to be steeper among women than men \citep{Khera2009SexCRPBodyFat}. These
observations provide biological and epidemiologic context for the present
finding, but they do not establish the mechanism underlying the sex
difference in HOMA-IR--hsCRP coupling observed here.

Adiposity appears to be an important component of this relationship. Adipose
tissue is metabolically and immunologically active, with adipose-derived
signals contributing to both insulin sensitivity and systemic inflammatory
biology \citep{Ouchi2011Adipokines}. In the present analysis, omission of
adiposity adjustment substantially increased both sex-specific HOMA-IR
slopes and the estimated sex contrast, whereas adjustment using BMI, waist
circumference, or waist-to-height ratio attenuated the contrast but did not
eliminate it. This pattern is consistent with adiposity accounting
statistically for a meaningful portion of the observed HOMA-IR--hsCRP
relationship, while the persistence of the sex contrast across several
anthropometric specifications suggests that the finding is not attributable
to any single representation of adiposity. Prior population-based work has
similarly reported a steeper association between total fat mass and CRP
among women than men \citep{Khera2009SexCRPBodyFat}. Because anthropometric
measures are imperfect proxies for body composition and fat distribution,
however, these analyses should not be interpreted as demonstrating an
association independent of adipose biology.

The estimated sex contrast was also largely unchanged across several
targeted sensitivity analyses addressing smoking, statin use, renal
function, hepatic status, alcohol use, and pregnancy. These factors were
evaluated because each could plausibly influence metabolic or inflammatory
measurements or their interpretation. Their inclusion or exclusion produced
little change in the estimated HOMA-IR--hsCRP sex interaction. These analyses
therefore reduce concern that the primary finding is readily explained by
any one of these measured factors, but they do not eliminate residual
confounding or establish that these factors are biologically irrelevant to
the observed association. For alcohol, comparison with a reference model
fitted to the same classifiable sample further indicated that the apparent
difference from the full-cohort estimate reflected sample restriction rather
than alcohol adjustment itself.

The most important qualification to the primary finding emerged when the
analysis was restricted to participants with hsCRP $\leq10$ mg/L. The
estimated sex contrast was attenuated by approximately 51\%, from $-0.0547$
in the full analytic population to $-0.0269$ after restriction, and the
confidence interval for the restricted estimate included zero. The
10-mg/L threshold has commonly been used in epidemiologic studies to reduce
the potential influence of acute inflammatory states when CRP is used as a
marker of chronic low-grade inflammation. However, values above this
threshold cannot be assumed to represent acute infection. Repeated-measure
studies have demonstrated that CRP concentrations above 10 mg/L may persist,
particularly among women and individuals with obesity, suggesting that
exclusion of all such observations can also remove participants with
sustained inflammatory phenotypes \citep{Ishii2012CRP}. Accordingly, the
attenuation observed here should not be interpreted simply as elimination
of acute inflammatory outliers. Rather, it indicates that the estimated sex
difference in HOMA-IR--hsCRP coupling is substantially influenced by
participants in the upper range of the hsCRP distribution.

Descriptive evaluation of participants with hsCRP $>10$ mg/L provided
additional context for this sensitivity. In both sexes, participants in the
upper hsCRP range had higher typical leukocyte counts, absolute neutrophil
counts, neutrophil percentages, and neutrophil-to-lymphocyte ratios than
participants with hsCRP $\leq10$ mg/L, consistent with greater inflammatory
or immune activation at the group level. The distribution of HOMA-IR,
however, differed more distinctly among women: median HOMA-IR was
substantially higher among women with hsCRP $>10$ mg/L than among women
below the threshold, whereas the corresponding male medians were nearly
identical. This pattern helps explain the geometry of the sensitivity
result: observations in the upper inflammatory range among women were also
shifted toward higher HOMA-IR, while the analogous male observations showed
less evidence of a comparable metabolic shift. These descriptive findings
do not establish a biological mechanism, but they indicate that the
attenuation after hsCRP restriction reflects structure in the joint
HOMA-IR--hsCRP distribution rather than merely the removal of a small number
of arbitrarily defined extreme observations.

Reproductive physiology represents an important potential direction for
subsequent investigation. Sex differences in insulin sensitivity vary
across the life course, and the relative metabolic advantage observed among
women before menopause appears to diminish after the menopausal transition;
sex hormones, adipose distribution, and inflammatory signaling have all
been proposed as contributing mechanisms
\citep{Gado2024SexIR,Henstridge2019MetabolicControlSex}. However,
reproductive state could not be classified reliably from age or recent
menstruation alone in the present NHANES sample. Hysterectomy, bilateral
oophorectomy, natural menopause, and exogenous hormone use represent
physiologically distinct states that would require explicit phenotype
definitions and separate consideration. We therefore did not perform post
hoc outcome modeling by menopausal or reproductive category. Whether the
observed sex difference in HOMA-IR--hsCRP coupling varies across reproductive
states remains an important question for future study.

Several limitations should be considered when interpreting these findings.
First, the cross-sectional design precludes determination of temporal or
causal direction between insulin resistance and systemic inflammation.
Second, HOMA-IR is a surrogate derived from fasting glucose and insulin
rather than a direct measure of insulin sensitivity such as the
hyperinsulinemic-euglycemic clamp
\citep{Matthews1985HOMA,Wallace2004HOMA}. This distinction may be
particularly relevant to the present study because relationships between
fasting glucose--insulin indices and clamp-measured insulin sensitivity can
vary by sex and other population characteristics
\citep{Pisprasert2013SurrogateIndices}. Accordingly, a one-unit difference
in HOMA-IR should not necessarily be assumed to represent an identical
physiological difference in insulin sensitivity in women and men. Third,
hsCRP is a nonspecific marker of systemic inflammation and does not identify
the source or chronicity of the inflammatory process
\citep{PepysHirschfield2003CRP}. Because hsCRP was measured at a single
examination, transient and persistent elevations could not be distinguished.
This limitation is particularly relevant given the substantial attenuation
of the estimated sex contrast after restriction to hsCRP $\leq10$ mg/L.

Selection and residual confounding also warrant consideration. The analytic
population was deliberately restricted to adults without diagnosed diabetes
and with HbA1c below the conventional prediabetes threshold, and the findings
therefore should not be assumed to generalize to individuals with
prediabetes or diabetes or to populations outside the NHANES target
population. Analyses additionally depended on participation in the fasting
laboratory subsample and complete availability of primary-model variables.
Appropriate fasting-subsample survey weights were used to support population
inference, and only 103 of 3,182 otherwise eligible participants were
excluded for incomplete primary-model covariates; nevertheless, selection
related to measurement availability or complete-case inclusion cannot be
excluded. Residual confounding also remains possible despite adjustment for
major metabolic covariates and the targeted sensitivity analyses. Race and
ethnicity, socioeconomic circumstances, physical activity, diet,
reproductive and hormonal factors, additional medication use, and more
detailed measures of body composition are among factors not comprehensively
represented in the primary model. Their omission should not be interpreted
as evidence that they are unimportant, but indiscriminate post hoc
adjustment for additional variables would not necessarily reduce bias and
could complicate interpretation.

The present analysis should also be viewed as a focused follow-up to an
exploratory observation from an earlier analysis of the same NHANES
population rather than as an independent replication. Sex-related
heterogeneity was identified during methodological evaluation of the
preceding study and motivated the specific HOMA-IR-by-sex interaction
evaluated here. The present study subsequently defined that interaction as
its primary question and subjected it to structured sensitivity analyses;
however, because the hypothesis was generated and evaluated using
overlapping data, the primary interaction estimate should not be interpreted
as independent confirmation of the original exploratory signal. Replication
in an independent population will therefore be important.

The study also has several strengths. The analytic population was
deliberately defined to examine metabolic--inflammatory heterogeneity among
adults who had not crossed conventional glycemic thresholds for prediabetes
or diagnosed diabetes. Sex differences were evaluated using a formal
interaction test rather than inferred from comparisons of statistical
significance within sex-specific subgroups. Analyses incorporated the
complex NHANES sampling design, fasting-subsample weights, and domain-based
variance estimation to support population-level inference. The primary
interaction was further evaluated across a structured series of sensitivity
analyses addressing exposure scale and influence, alternative
representations of adiposity, several clinical and behavioral factors,
pregnancy, and the upper range of the hsCRP distribution. Where sensitivity
analyses altered the available analytic sample, same-sample reference models
were used to distinguish changes associated with covariate adjustment from
those associated with sample restriction. Finally, sensitivity findings
were evaluated on the basis of changes in effect estimates and uncertainty
rather than solely according to whether a nominal threshold for statistical
significance was crossed.

In conclusion, among U.S. adults without diagnosed diabetes and with HbA1c
below the conventional threshold for prediabetes, the adjusted association
between HOMA-IR and hsCRP was more positive among women than men. This sex
difference persisted across alternative HOMA-IR scales, progressively
trimmed exposure distributions, multiple representations of adiposity, and
a range of targeted clinical and behavioral sensitivity analyses. However,
restriction to hsCRP $\leq10$ mg/L reduced the estimated sex contrast by
approximately one-half, indicating that the upper portion of the
inflammatory distribution contributes importantly to the observed
heterogeneity. Taken together, these findings suggest that sex-related
differences in metabolic--inflammatory coupling may be detectable even
below conventional glycemic thresholds, while also demonstrating that the
magnitude of this difference depends substantially on the inflammatory
range under study. Independent replication and studies incorporating more
direct measures of insulin sensitivity, repeated inflammatory measurements,
body composition, and reproductive physiology will be important for
determining the biological basis and clinical relevance of this pattern.

% ============================================================

\section*{Author Contributions}

S.M. conceived the study, developed the analytic framework, conducted
the statistical analyses, interpreted the results, and drafted the
manuscript. J.B. contributed to study development, interpretation of
the findings, and critical revision of the manuscript. Both authors
reviewed and approved the final manuscript.

\section*{Funding}

This research received no external funding.

\section*{Conflicts of Interest}

The authors declare no conflicts of interest.

\section*{Ethics Statement}

This study used publicly available, de-identified data from the National
Health and Nutrition Examination Survey (NHANES). NHANES protocols were
approved by the National Center for Health Statistics Research Ethics Review
Board, and written informed consent was obtained from participants in
accordance with NHANES procedures. Because the present study involved
secondary analysis of publicly available, de-identified data, no additional
institutional review board approval or informed consent was required.

\section*{Data Availability}

The NHANES data used in this study are publicly available from the National
Center for Health Statistics.

% ============================================================
% References
% ============================================================

\bibliographystyle{plainnat}
\bibliography{manuscript_references}

\end{document}
